\documentclass[UTF8,fontset=fandol,12pt,a4paper]{ctexart}
% Individual identity supplied by the author on 2026-11-09.
\newcommand{\AuthorName}{蒋佳豪}
\newcommand{\StudentID}{22671064}
\newcommand{\CourseName}{工程技术前沿}
\newcommand{\ReviewTitle}{面向科研任务的多智能体系统综述}
\newcommand{\ReviewSubtitle}{知识、记忆与协作的机制及证据边界}
\newcommand{\ReportDate}{2026年11月9日}

\usepackage{review}
\begin{document}
\reviewcover{中文学术综述排版样稿（非提交稿）}
\section*{样稿说明}
本文件只展示封面、章节、图表和参考文献的排版。它不是已完成的课程作业，
不满足正文不少于3000字的要求。课程截图明确禁止直接提交AI生成内容；
下列写作提示应由作者在阅读原文后，用自己的表述和判断替换。
作者姓名与学号尚未确认，不能按正式文件名打包提交。

\section*{摘要}
\writingtask{正文完成后撰写约200至300字。交代研究问题、文献范围、方法分类和你的主要判断。
不要仅写“多智能体很重要”，也不要把拟议方案写成已完成的实验。}
\noindent\textbf{关键词：}多智能体；科研辅助；检索增强生成；记忆机制；可信性

\section{问题引入与范围界定}
\writingtask{约550字。从一次具体科研任务切入：检索文献、核验观点、设计实验或复现代码。
解释单次问答为何不足，再界定本文讨论的是基于大语言模型的任务协作系统，
不泛化到所有多智能体强化学习。说明本题与课堂知识增强、记忆和科研验证的联系。}

\subsection{研究问题}
\writingtask{写出一个贯穿全文的问题，例如：协作系统如何在长任务中维持可追溯、可更新的证据？
“增加角色数量是否等于提高可靠性”可作比较问题，但答案需要文献或推理支持。}

\subsection{文献选择}
\writingtask{交代检索关键词、检索日期和纳入标准。当前8篇是有目的的代表性选取，
不应称为穷尽式或系统性综述；如需声称系统综述，应补充可复现筛选过程。}

\clearpage
\section{方法分类与工作机制}
\writingtask{约1300字。按机制组织，不要每篇论文各写一段却缺少比较。
可依次说明外部知识、工作记忆、角色协作与工具反馈之间的关系。}
\subsection{外部证据：从检索增强到图结构组织}
\writingtask{阅读RAG与GraphRAG\cite{rag,graphrag}。区分局部事实检索、跨文档关系和语料级总结，
再解释图索引可能增加哪些构建成本与抽取误差。不能据此声称所有图检索都优于向量检索。}
\subsection{任务状态：记忆检索与上下文管理}
\writingtask{阅读长上下文、Generative Agents及MemGPT\cite{middle,generative,memgpt}。
区分外部事实来源与代理自身经历。讨论错误记忆、过期信息及信息位置的影响，
不要把“能记住”直接等同于“记住的事实正确”。}
\subsection{协作控制：角色划分、行动和反馈}
\writingtask{阅读ReAct、AutoGen及Agent Laboratory\cite{react,autogen,lab}。
解释单代理工具交互与多代理协作并不相同，说明人类反馈在何处进入流程。
比较系统完成了哪些任务、用了哪些评估，避免把自动写报告理解为自动发现真理。}

\begin{table}[htbp]\centering\small
\caption{作者填写的文献比较框架；空白不表示零表现。}
\begin{tabularx}{\linewidth}{@{}lXXX@{}}\toprule
方法类别 & 输入与状态 & 主要收益 & 证据与局限\\\midrule
知识增强 & 待读文献后填写 & 待读文献后填写 & 待读文献后填写\\
记忆管理 & 待读文献后填写 & 待读文献后填写 & 待读文献后填写\\
多代理协作 & 待读文献后填写 & 待读文献后填写 & 待读文献后填写\\\bottomrule
\end{tabularx}\end{table}

\clearpage
\section{可信性评估与方法比较}
\writingtask{约650字。选择共同维度，例如引用是否支持断言、代码能否复现、
实验条件是否公平、人类修改成本和任务失败率。讨论不同论文的任务、模型及预算差异，
不要把不同测试集上的准确率直接排名，也不要虚构本作业没有运行的实验。}

\section{个人理解与讨论}
\writingtask{约850字。本节必须体现你自己的阅读和判断。选一个课堂提到的科研例子，
先写支持证据，再说明反例或适用边界，最后给出你的取舍。
可以判断何时需要图结构、何时简单检索已经足够，以及哪些决定应由研究者保留。}
\subsection{你赞同或质疑什么}
\writingtask{写下一个明确观点，并指出哪篇论文的哪一节支持它。
对“多个代理能互相验证”这一说法，追问：它们是否共享同一个错误来源？}
\subsection{一个可以检验的改进设想}
\writingtask{约定输入、对照方法、指标及失败条件。若没有真正做实验，明确标为研究设想，
不要出现“实验表明”或虚构数值。}

\section{结论}
\writingtask{约250字。回到开头问题，给出条件化判断，指出还无法回答的部分。
建议正文目标约3600字，最终以排除封面、摘要、参考文献和提示文字后的实际正文核算。}

\clearpage
\begingroup\small\raggedright
\begin{thebibliography}{99}
\bibitem{rag} Lewis P, Perez E, Piktus A, et al. Retrieval-Augmented Generation for Knowledge-Intensive NLP Tasks[EB/OL]. 2020. arXiv:2005.11401v4. \url{https://arxiv.org/abs/2005.11401v4}.
\bibitem{graphrag} Edge D, Trinh H, Cheng N, et al. From Local to Global: A Graph RAG Approach to Query-Focused Summarization[EB/OL]. 2024. arXiv:2404.16130v2. \url{https://arxiv.org/abs/2404.16130v2}.
\bibitem{middle} Liu N F, Lin K, Hewitt J, et al. Lost in the Middle: How Language Models Use Long Contexts[EB/OL]. 2023. arXiv:2307.03172v3. \url{https://arxiv.org/abs/2307.03172v3}.
\bibitem{react} Yao S, Zhao J, Yu D, et al. ReAct: Synergizing Reasoning and Acting in Language Models[EB/OL]. 2022. arXiv:2210.03629v3. \url{https://arxiv.org/abs/2210.03629v3}.
\bibitem{autogen} Wu Q, Bansal G, Zhang J, et al. AutoGen: Enabling Next-Gen LLM Applications via Multi-Agent Conversation[EB/OL]. 2023. arXiv:2308.08155v2. \url{https://arxiv.org/abs/2308.08155v2}.
\bibitem{generative} Park J S, O'Brien J C, Cai C J, et al. Generative Agents: Interactive Simulacra of Human Behavior[EB/OL]. 2023. arXiv:2304.03442v2. \url{https://arxiv.org/abs/2304.03442v2}.
\bibitem{memgpt} Packer C, Wooders S, Lin K, et al. MemGPT: Towards LLMs as Operating Systems[EB/OL]. 2023. arXiv:2310.08560v2. \url{https://arxiv.org/abs/2310.08560v2}.
\bibitem{lab} Schmidgall S, Su Y, Wang Z, et al. Agent Laboratory: Using LLM Agents as Research Assistants[EB/OL]. 2025. arXiv:2501.04227v2. \url{https://arxiv.org/abs/2501.04227v2}.
\end{thebibliography}
\endgroup

\end{document}
